\documentclass[header={Ancillary Note: Basic Topology},notheorems]{study-note}

\begin{document}

\StudyTitleBlock{N3}{Ancillary Note}{Dense, open, closed, compact, and related topological concepts}

\vspace{0.8em}

\begin{focusbox}
These words appear constantly in analysis. This note is a quick reference for the basic topological vocabulary used throughout functional analysis.
\end{focusbox}

\tableofcontents

\section{Open And Closed Sets}

\begin{defbox}{Open set}
A subset $U$ of a metric space $X$ is \textbf{open} if for every $x\in U$ there exists $r>0$ such that
\[
B(x,r)\subset U.
\]
\end{defbox}

\begin{defbox}{Closed set}
A subset $F\subset X$ is \textbf{closed} if its complement $X\setminus F$ is open.
Equivalently, every limit of a convergent sequence in $F$ still belongs to $F$.
\end{defbox}

\begin{warningbox}{Common pitfall}
A set can be neither open nor closed, and it can also be both open and closed. In a connected space, the only clopen sets are the empty set and the whole space.
\end{warningbox}

\section{Closure, Interior, Boundary}

\begin{defbox}{Closure}
The \textbf{closure} of $A\subset X$ is the smallest closed set containing $A$, denoted by $\ol A$.
\end{defbox}

\begin{defbox}{Interior}
The \textbf{interior} of $A$ is the largest open set contained in $A$, denoted by $A^{\circ}$.
\end{defbox}

\begin{defbox}{Boundary}
The \textbf{boundary} of $A$ is
\[
\partial A:=\ol A\setminus A^{\circ}.
\]
\end{defbox}

\section{Dense Sets}

\begin{defbox}{Dense set}
A subset $A\subset X$ is \textbf{dense} in $X$ if
\[
\ol A=X.
\]
\end{defbox}

\begin{thmbox}{Equivalent viewpoint}
$A$ is dense in $X$ if and only if every nonempty open set in $X$ intersects $A$.
\end{thmbox}

\begin{exbox}{Standard examples}
\begin{itemize}
\item The rational numbers are dense in $\R$.
\item Polynomials are dense in many function spaces, but not equal to them.
\item Finite linear combinations of basis vectors are dense in $\ell^2$, but they are not all of $\ell^2$.
\end{itemize}
\end{exbox}

\section{Compactness}

\begin{defbox}{Compact set}
A subset $K$ of a metric space is \textbf{compact} if every open cover has a finite subcover.
\end{defbox}

\begin{thmbox}{Sequential form in metric spaces}
For metric spaces, $K$ is compact if and only if every sequence in $K$ has a convergent subsequence whose limit lies in $K$.
\end{thmbox}

\begin{thmbox}{Compactness in metric spaces}
For a subset $K$ of a metric space, the following are equivalent:
\begin{enumerate}
\item $K$ is compact;
\item $K$ is sequentially compact;
\item $K$ is complete and totally bounded.
\end{enumerate}
\end{thmbox}

\begin{warningbox}{Infinite-dimensional warning}
In infinite-dimensional normed spaces, closed and bounded does \emph{not} imply compact. This is one of the main differences from finite-dimensional Euclidean spaces.
\end{warningbox}

\section{Nowhere Dense And Meager}

\begin{defbox}{Nowhere dense}
A set $A\subset X$ is \textbf{nowhere dense} if the interior of $\ol A$ is empty.
\end{defbox}

\begin{defbox}{Meager set}
A set is \textbf{meager} if it is a countable union of nowhere dense sets.
\end{defbox}

\begin{focusbox}
These notions are important because Baire's theorem says a complete metric space is not meager in itself.
\end{focusbox}

\section{What To Remember Quickly}

\begin{focusbox}
Quick summary:
\begin{enumerate}
\item open = contains a little ball around each point;
\item closed = contains all its limit points;
\item dense = closure is everything;
\item compact in metric spaces = every sequence has a convergent subsequence;
\item in infinite-dimensional spaces, closed and bounded is usually not compact.
\end{enumerate}
\end{focusbox}

\end{document}
